本章把第~\ref{sec:props} 节的理论结论放到原文实证节（第 424--650 行）中逐项对照。理论刻画的是“风险集已知、一代传播、报告率已知”条件下估计量的方差结构。实证使用的是由报告病例构造的周度增长代理，两者之间隔着若干步近似。本章先说明数据和变量的构造，再确定理想条件下仿射模型参数对应的理论量，并解释实际估计值偏离这些理论量的原因。随后讨论 FGLS、聚类 Bootstrap 与工具变量成立的条件，最后逐张解读原文的实证表。每一处偏离都给出定量依据，并标注第~\ref{sec:checks} 节问题表中的编号（I-xx）。

\subsection{数据构造与仿射模型}

原文使用全美各州及华盛顿特区医院向 CDC NHSN 上报的周度实验室确诊流感住院人数与各州医院报告覆盖率 \citep{cdc2026nhsn}。该住院指标自 2021--22 赛季起是 CDC FluSight 的核心预测目标 \citep{mathis2024evaluation}。样本覆盖 2020 年 8 月至 2026 年 9 月、50 个州及华盛顿特区，共 16,218 个州–周观测。其中 504 个州–周的住院数缺失，原文按 0 例处理（第一阶段 468 个、第二阶段 36 个）。制度切换期的“未填报”与“零住院”含义不同，原文第 430 行已把这一处理的影响限定在第二阶段。

样本筛选只有一条规则：历史有效输入压力 $D_t\ge5$ 的州–周进入推断样本，共 10,303 个；原文不另设覆盖率门槛，以免系统性剔除低覆盖观测。这 10,303 个观测按报告制度分为三个阶段，见表~\ref{tab:em-phase}。第二阶段覆盖率从 91.1\% 跌到 54.6\%，原文把它当作报告率杠杆 $\rho$ 的制度层面代理，第~\ref{sec:em-space} 小节检验这一代理能解释多少观测噪声的变化。

\begin{table}[!tbp]
\centering
\small
\caption{原文三个报告制度阶段的划分与覆盖率（依据原文第 430--437 行）}
\label{tab:em-phase}
\begin{tabular}{L{2.6cm}L{4.4cm}L{1.1cm}L{1.3cm}L{1.9cm}L{1.7cm}}
\toprule
阶段 & 时间范围 & 有效周 & 样本数 $N$ & 覆盖率（样本平均） & 覆盖率中位数 \\
\midrule
前过渡高覆盖期 & 2020 年 10 月至 2024 年 4 月 30 日 & 184 & 6,009 & 91.1\% & 92.3\% \\
自愿过渡期 & 2024 年 5 月 1 日至 2024 年 10 月 31 日 & 26 & 420 & 54.6\% & 45.7\% \\
CMS 新规重制期 & 2024 年 11 月 1 日至 2026 年 9 月初 & 96 & 3,874 & 91.8\% & 93.2\% \\
\bottomrule
\end{tabular}
\par\smallskip
\begin{minipage}{0.92\linewidth}\footnotesize
注：样本数为 $D_t\ge5$ 筛选后的州–周数，三阶段合计 10,303。第一阶段合并了大流行早期段（$N=2{,}627$，覆盖率 87.8\%）与 2022 年 10 月起的法定强制段（$N=3{,}382$，覆盖率 93.7\%）。
\end{minipage}
\end{table}

原文用滞后 1--3 周报告病例的加权和定义第 $j$ 州第 $t$ 周的历史有效输入压力，构造如下式所示：
\begin{equation*}
D_{j,t} = \sum_{s=1}^3 w_s C_{j, t+1-s}.
\tag{原文式 30}
\end{equation*}
主分析取 $w=(0.65,0.25,0.10)$，平均回溯延迟为 $\sum_s s\,w_s=1\times0.65+2\times0.25+3\times0.10=1.45$ 周，乘以 7 得 10.15 天，原文写作“约 10.2 天”。增长代理定义为 $G_{j,t+1}=C_{j,t+1}/D_{j,t}$。敏感性分析另取三周等权 $(1/3,1/3,1/3)$ 与单周滞后 $(1,0,0)$，平均回溯延迟分别为 2.0 周与 1.0 周。后文反复用到权重的平方和 $\sum_s w_s^2$：主权重为 $0.65^2+0.25^2+0.10^2=0.4225+0.0625+0.0100=0.495$，单周为 1，三周等权为 $1/3$。

主分析只用流行初期上升波段：2022--23 季的 10 月 8 日至 11 月 26 日、2023--24 季的 10 月 14 日至 12 月 16 日、2024--25 季的 11 月 9 日至 1 月 18 日。波段内有效样本 $N=1267$，阶段平均增长比 $\bar R_{\mathrm{early}}=1.6706$（标准差 0.7116）。原文把命题 2 的 $v_R+A/m$ 改写成相对尺度后，对原始观测估计如下的仿射条件期望模型：
\begin{equation*}
\mathbb{E}\!\left[\left(\frac{G_{t+1} - \bar{R}_{\mathrm{early}}}{\bar{R}_{\mathrm{early}}}\right)^{2}\ \middle|\ D_t\right] = a + \frac{b}{D_t},
\tag{原文式 31}
\end{equation*}
左端的被解释变量记为 $L=\big((G-\bar R_{\mathrm{early}})/\bar R_{\mathrm{early}}\big)^2$。由原文式 27 的逐观测恒等式 $\big((C-\bar RD)/(\bar RD)\big)^2=\big((G-\bar R)/\bar R\big)^2$，$L$ 同时是事前单步预测 $\widehat C_{t+1\mid t}=\bar R_{\mathrm{early}}D_t$ 的相对平方误差，所以原文式 31 可以从“增长率估计”与“病例预测”两个角度解读。表~\ref{tab:em-mapping} 把理论量与实证量逐一对齐，并列出每一处对齐所需的假设；后续各小节分别处理这些假设失效的后果。

\begin{table}[!tbp]
\centering
\small
\caption{理论量与实证量的对应关系及各自需要的假设}
\label{tab:em-mapping}
\begin{tabular}{L{3.2cm}L{3.6cm}L{6.6cm}}
\toprule
理论量 & 实证对应量 & 对应成立需要的条件与实际偏离 \\
\midrule
期望报告规模 $m=\rho n$ & $D_{j,t}$（报告病例加权和） & 需要 $D=\rho\Lambda$ 无噪；实际 $D$ 含二项报告噪声，抬高 $b$（第~\ref{sec:em-delta} 小节，I-13） \\
估计量 $C_t/(\rho n)$ & $G_{j,t+1}=C_{j,t+1}/D_{j,t}$ & 需要 $G$ 对应一代传播；实际周增长比跨越约 2 个以上世代，抬高 $a$（第~\ref{sec:em-intercept} 小节，I-12） \\
跨期水平 $\bar R$ & 阶段平均 $\bar R_{\mathrm{early}}=1.6706$ & 需要基准无偏且外生；实际为同一样本的回顾性估计，事前滚动基准下误差整体抬升 \\
条件矩 $\mathbb{E}[R\mid n]=\bar R$ & $D_t$ 外生 & 时间序列中 $D_t$ 由过去传播决定，可能与 $R_t$ 相关（I-04） \\
$v_R/\bar R^2$ 与 $A/\bar R^2$ & 截距 $a$ 与斜率 $b$ & 仅在上述条件全部成立时 $a$、$b$ 分别等于这两个量，$b/a=m_\times$ \\
\bottomrule
\end{tabular}
\end{table}

\subsection{理想映射与量级比较}

\subsubsection{理想条件下 $a$、$b$ 与 $b/a$ 的含义}

原文没有写出 $a$、$b$ 与理论参数的精确对应，只在第 508 行把 $a$ 与 $v_R/\bar R^2$ 作量级比较。下面补齐理想映射。设四个条件成立：(i) $D=\rho\Lambda$ 精确，分母无噪；(ii) $G$ 对应一代的再生数；(iii) $\mathbb{E}[R\mid D]=\bar R$，$\operatorname{Var}(R\mid D)=v_R$；(iv) 基准 $\bar R$ 无偏。
\begin{enumerate}
\item 在条件 (i) 下 $G=C/(\rho\Lambda)$，它恰好是命题 2 的估计量 $C_t/(\rho n)$，期望报告规模 $m$ 就是 $D$。
\item 由塔式性质与条件 (iii) 得 $\mathbb{E}[G\mid D]=\mathbb{E}\big[\mathbb{E}[G\mid D,R]\mid D\big]=\mathbb{E}[R\mid D]=\bar R$，再由条件 (iv)，$G-\bar R$ 的条件期望为零，于是 $\mathbb{E}[(G-\bar R)^2\mid D]=\operatorname{Var}(G\mid D)$。
\item 由全方差公式（附录引理~\ref{lem:ap-totvar}）与命题 2 的结论得 $\operatorname{Var}(G\mid D)=v_R+A/D$，其中 $A=\bar R+\rho(\bar R^2+v_R)/k$ 为原文式 8 定义的规模相关方差系数。
\item 两边除以常数 $\bar R^2$ 即得式~\ref{eq:em-ideal}，与原文式 31 逐项比对可读出 $a=v_R/\bar R^2$、$b=A/\bar R^2$，于是 $b/a=A/v_R=m_\times$，它与 $\bar R$ 无关。
\end{enumerate}
\begin{equation}\label{eq:em-ideal}
\mathbb{E}\Big[\Big(\frac{G-\bar R}{\bar R}\Big)^2\,\Big|\,D\Big]=\frac{\operatorname{Var}(G\mid D)}{\bar R^2}=\frac{v_R}{\bar R^2}+\frac{A/\bar R^2}{D}.
\end{equation}
蒙特卡洛核验在条件 (i)--(iv) 下取 $\bar R=1.67$、$v_R=0.06$、$k=0.35$、$\rho=0.021$（$A=1.8409$），FGLS 截距为 0.0216（SE 0.0005），理论值 $v_R/\bar R^2=0.0215$；斜率为 0.6806（SE 0.0153），理论值 $A/\bar R^2=0.6601$。理想映射本身没有问题，问题出在实证数据不满足其中的条件 (i)--(iv)。

\subsubsection{情景隐含值与实证估计的比较（I-14）}

原文第 515 行给出流感住院情景：$\bar R_{\mathrm{early}}\approx1.67$，超传播参数 $k\in[0.2,0.5]$ \citep{lloydsmith2005superspreading}，住院比例 $\rho\in[0.0056,0.021]$（由 \citet{zhou2012hospitalizations} 的住院负担与 \citet{tokars2018seasonal} 的侵袭率折算），$v_R\in[0.04,0.08]$。在给出 $m_\times$ 两端的角点上计算：
\begin{align*}
A_{\min}&=1.67+0.0056\times2.8689/0.5=1.7021, & A_{\max}&=1.67+0.021\times2.8289/0.2=1.9670,\\
m_\times&\in[1.7021/0.08,\ 1.9670/0.04]=[21.28,\ 49.18]. &&
\end{align*}
同一情景隐含 $a=v_R/\bar R^2\in[0.04/2.7889,\ 0.08/2.7889]=[0.0143,\ 0.0287]$，$b=A/\bar R^2\in[0.610,\ 0.705]$；若在全部角点上取极值，$b$ 的上端为 0.707，下文倍数结论不受影响。原文把 $m_\times$ 上端写作 49.3，来自 $A$ 中取 $v_R=0.08$、分母取 0.04 的混用（I-08）。表~\ref{tab:em-compare} 把情景隐含值与实证 FGLS 估计并排列出。

\begin{table}[!tbp]
\centering
\small
\caption{仿射模型参数的情景隐含值与实证 FGLS 估计}
\label{tab:em-compare}
\begin{tabular}{L{2.6cm}L{3.2cm}L{4.4cm}L{3.2cm}}
\toprule
参数 & 情景隐含范围 & 实证估计（州聚类 95\% 区间） & 实证与隐含之比 \\
\midrule
截距 $a$ & $0.0143$--$0.0287$ & $0.0991$ $[0.0841, 0.1173]$ & 3.5--6.9 倍 \\
斜率 $b$ & $0.610$--$0.705$ & $2.2434$ $[1.5019, 3.0161]$ & 3.2--3.7 倍 \\
$b/a$ 与 $m_\times$ & $21.28$--$49.18$ & $22.6$ $[13.7, 33.7]$ & 落在区间内 \\
\bottomrule
\end{tabular}
\end{table}

原文第 508 行写道“经验值 0.0991 与参数空间隐含值 0.014--0.029 相符”。按表~\ref{tab:em-compare}，$a$ 是隐含值的 3.5--6.9 倍，“相符”不成立；$b$ 同样被放大到 3.2--3.7 倍。$b/a=22.6$ 能落在 21.3--49.3 之内，原因是 $a$ 与 $b$ 被放大了相近的倍数：$b$ 的放大主要来自随机分母（第~\ref{sec:em-delta} 小节），$a$ 的放大主要来自周增长比与世代再生数的口径差（第~\ref{sec:em-intercept} 小节）。两种机制互不相关，倍数接近只是巧合。分母噪声单独作用时，蒙特卡洛给出 $b/a=59.2$，而真值 $A/v_R=30.7$，比值会偏离一倍。所以 $b/a$ 与情景区间同量级不能当作结构验证。原文第 508 行宜改为“截距高于情景隐含值（3.5--6.9 倍），差额来自 $G$ 与世代 $R$ 的口径差及代理偏差”，第 515 行宜改为“$b/a$ 与情景区间同量级，但 $a$、$b$ 各自均高于隐含值，这一吻合不构成结构识别”。

\subsubsection{原文表 2 的读法}

原文表 2（第 483--501 行）把初期波段 1267 个观测按 $D_t$ 分成六个箱，表~\ref{tab:em-scale} 摘录了其中的关键列。第三列“相对均方误差”是损失 $L$ 的箱内样本平均，它是原文式 31 左端条件期望的经验对应物；理想映射下它对应 $(v_R+A/m)/\bar R^2$ 在箱内的平均，即式~\ref{eq:em-ideal}。六个箱的数值按样本数加权平均等于全体的 0.1813，核验在四位小数上复现了这一点，它也是第~\ref{sec:em-budget} 小节误差预算的分母。

\begin{table}[!tbp]
\centering
\small
\caption{原文表 2 摘录：初期波段单步相对平方误差按规模分箱的统计量}
\label{tab:em-scale}
\setlength{\tabcolsep}{4pt}
\begin{tabular}{L{2.2cm}L{1.4cm}L{2.0cm}L{1.6cm}L{1.6cm}L{1.6cm}L{3.0cm}}
\toprule
规模箱 $D_t$ & 样本数 $N$ & 相对均方误差 & 中位数 & 普通 SEM & 州聚类 SE & 州聚类 95\% 区间 \\
\midrule
$[5, 20)$ & 319 & 0.3292 & 0.0969 & 0.0425 & 0.0437 & $[0.2435, 0.4149]$ \\
$[20, 50)$ & 297 & 0.1605 & 0.0596 & 0.0169 & 0.0189 & $[0.1235, 0.1974]$ \\
$[50, 100)$ & 187 & 0.1446 & 0.0483 & 0.0194 & 0.0198 & $[0.1057, 0.1835]$ \\
$[100, 250)$ & 220 & 0.1118 & 0.0545 & 0.0110 & 0.0114 & $[0.0895, 0.1341]$ \\
$[250, 600)$ & 152 & 0.1032 & 0.0460 & 0.0097 & 0.0133 & $[0.0771, 0.1293]$ \\
$[600, 5000)$ & 92 & 0.1054 & 0.0735 & 0.0109 & 0.0146 & $[0.0769, 0.1340]$ \\
\bottomrule
\end{tabular}
\end{table}

读这张表要抓三点。其一，中位数远低于箱内样本平均，$[5,20)$ 箱为 0.0969 对 0.3292，说明损失右偏重尾；这正是小规模箱 $\operatorname{Var}(L)/\mathbb{E}[L]^2$ 高达 5.32 的原因，也解释了剔除损失最大的 1\% 观测后 $b/a$ 从 22.6 降到 10.6。其二，州聚类 SE 普遍大于普通 SEM，$[250,600)$ 箱为 0.0133 对 0.0097，反映同州跨周观测的正相关。其三，第 511 行的“下降 51.3\%”即 $(0.3292-0.1605)/0.3292=0.5125$；$[5,20)$ 箱区间下界 0.2435 高于 $[20,50)$ 箱上界 0.1974，两箱在聚类口径下也分离。大规模三箱 0.1118、0.1032、0.1054 不再单调下降，与截距 $a=0.0991$ 形成的平台一致。

\subsection{随机分母：一阶 delta 法}\label{sec:em-delta}

理想映射的条件 (i) 要求分母无噪，而原文式 30 的 $D_{j,t}$ 由报告病例构造，本身带有报告抽样噪声。原文式 28 后的说明已提到“随机分母 $D_t$ 对条件方差的倒数二次方重新加权”，局限性第 664 行也只作定性叙述。下面用一阶 delta 法把这一效应写成可计算的形式（I-13）。设过去各周真实感染 $I_s\equiv\Lambda$，报告 $C_s\sim\mathrm{Bin}(\Lambda,\rho)$ 相互独立，$D=\sum_s w_sC_s$，并记 $D=\rho\Lambda(1+\eta)$。
\begin{enumerate}
\item 由二项分布的矩 $\mathbb{E}[C_s]=\rho\Lambda$、$\operatorname{Var}(C_s)=\rho(1-\rho)\Lambda$ 与 $\sum_s w_s=1$ 得 $\mathbb{E}[D]=\rho\Lambda$，所以 $\mathbb{E}[\eta]=0$；由各周独立得 $\operatorname{Var}(D)=\sum_s w_s^2\rho(1-\rho)\Lambda$，于是 $\operatorname{Var}(\eta)=(1-\rho)\sum_s w_s^2/(\rho\Lambda)=(1-\rho)\sum_s w_s^2/m$。
\item 记 $G_0=C/(\rho\Lambda)$，则 $G=C/D=G_0/(1+\eta)$；由一阶 Taylor 展开 $1/(1+\eta)\approx1-\eta$ 得 $G\approx G_0(1-\eta)$。给定 $\Lambda$，过去各周的报告抽样与未来一周的传播相互独立，所以 $G_0$ 与 $\eta$ 独立。
\item 对相互独立的 $X$ 与 $Y=1-\eta$，由 $\mathbb{E}[Y]=1$、$\mathbb{E}[Y^2]=1+\operatorname{Var}(\eta)$ 可得 $\operatorname{Var}(XY)=\mathbb{E}[X^2]\mathbb{E}[Y^2]-\mathbb{E}[X]^2=\operatorname{Var}(X)+\mathbb{E}[X^2]\operatorname{Var}(\eta)$。
\item 由命题 2 得 $\operatorname{Var}(G_0)=v_R+A/m$，由二阶矩关系得 $\mathbb{E}[G_0^2]=\bar R^2+v_R+O(1/m)$，把两者与第 1 步的 $\operatorname{Var}(\eta)$ 代入第 3 步的公式即得式~\ref{eq:em-delta}。
\item 均值方面，由 $\mathbb{E}[1/(1+\eta)]-1=O(1/m)$ 可知偏差平方是 $O(1/m^2)$，在 $a+b/D$ 的精度下可以忽略，所以分母噪声只改变斜率，截距 $a$ 保持 $v_R/\bar R^2$。
\end{enumerate}
\begin{equation}\label{eq:em-delta}
\operatorname{Var}(G\mid m)\approx v_R+\frac{A+(1-\rho)(\bar R^2+v_R)\sum_s w_s^2}{m},
\qquad
b\approx\frac{A+(1-\rho)(\bar R^2+v_R)\sum_s w_s^2}{\bar R^2}.
\end{equation}
若报告存在过离散，报告方差写作 $\varphi\rho I$（$\varphi=1$ 对应 Poisson 抽稀），同样的步骤给出
\begin{equation*}
b\approx\big[\varphi\bar R+\rho\bar R+\rho(\bar R^2+v_R)/k+\varphi(\bar R^2+v_R)\textstyle\sum_s w_s^2\big]/\bar R^2 .
\end{equation*}

取 $\bar R=1.67$、$v_R=0.06$、$k=0.35$、$\rho=0.021$，则 $A=1.8409$，$A/\bar R^2=0.6601$。主权重下式~\ref{eq:em-delta} 的第二项为 $0.979\times2.8489\times0.495/2.7889\approx0.4950$，合计 $b\approx1.1551$；单周与三周等权分别得 1.6602 与 0.9934。表~\ref{tab:em-delta} 把 delta 法预测、蒙特卡洛结果与原文表 1 的实证比值放在一起。

\begin{table}[!tbp]
\centering
\small
\caption{三种滞后权重下斜率 $b$ 的 delta 法预测、蒙特卡洛结果与实证比值}
\label{tab:em-delta}
\setlength{\tabcolsep}{4pt}
\begin{tabular}{L{2.4cm}L{1.1cm}L{1.4cm}L{1.6cm}L{2.2cm}L{1.6cm}L{1.6cm}L{1.4cm}}
\toprule
权重 & $\sum w^2$ & 平均滞后（周） & 预测 $b$ & 蒙特卡洛 $b$（SE） & 预测比值 & 蒙特卡洛比值 & 实证比值 \\
\midrule
主权重 $(0.65,0.25,0.10)$ & 0.495 & 1.45 & 1.155 & 1.195（0.0450） & 1 & 1 & 1 \\
单周 $(1,0,0)$ & 1 & 1.0 & 1.660 & 1.566（0.0336） & 1.437 & 1.311 & 1.397 \\
三周等权 & 1/3 & 2.0 & 0.993 & 0.999（0.0396） & 0.860 & 0.836 & 0.859 \\
\bottomrule
\end{tabular}
\par\smallskip
\begin{minipage}{0.92\linewidth}\footnotesize
注：比值均相对主权重。实证比值取原文表 1 的 $b$：$3.1337/2.2434=1.397$，$1.9262/2.2434=0.859$。无分母噪声时三种权重的 $b$ 都是 0.6601。
\end{minipage}
\end{table}

表~\ref{tab:em-delta} 支持四个结论。第一，三种权重下 $b$ 的排序与比值由 $\sum w^2$ 定量解释：预测比值 1.437 与 0.860 对实证 1.397 与 0.859，原文第 445 行所说“经验拐点对回溯窗口长度存在敏感性”由此有了机制来源。原文表 1 中三种权重的有效样本与 $\bar R_{\mathrm{early}}$ 各不相同，所以这里的比值比较是近似对照。第二，水平上二项分母噪声把 $b$ 从 0.6601（实证的 29\%）抬到 1.1551（实证的 51\%），剩余缺口可由 $\varphi\approx1.96$ 的报告过离散补足。$\varphi=1.96$ 时 delta 法预测 2.2386，蒙特卡洛为 2.0190（SE 0.0343），量级与实证 2.2434 接近，但截距仍停在 0.0227，$a$ 的放大需要另找来源。过离散只是一种可能解释，数据本身不能把它与其他机制区分开。第三，一阶 delta 法在分母噪声大时高估约 5\%--10\%，原因是含噪回归元 $1/D$ 的回归稀释与 $D\ge5$ 截断，单周权重下 1.660 对 1.566 即属此类。第四，分母噪声只抬高斜率，主权重下蒙特卡洛的 $b/a$ 被推到 59.2，而 $A/v_R=30.7$，所以分母含噪时 $b/a$ 会高估 $m_\times$。实证 $b$ 因此不能直接读作 $A/\bar R^2$。

\subsection{截距与口径差}\label{sec:em-intercept}

原文表 1 的截距列给出一组关键证据：单周权重（平均滞后 1.0 周）$a=0.0682$，主权重（1.45 周）$a=0.0991$，三周等权（2.0 周）$a=0.1524$，随平均滞后单调增大。若 $a$ 只是 $v_R/\bar R^2$，它应当与权重无关；第~\ref{sec:em-delta} 小节的蒙特卡洛也显示二项分母噪声不改变截距，三种权重下分别为 0.0225、0.0202、0.0214，都在理论值 0.0215 附近。截距随滞后增大必须由分母噪声之外的机制解释（I-12）。

最直接的解释是口径差。理论中的 $R_t$ 是单代再生数，流感的世代间隔约 3 天，一周约跨 2 个世代。$G=C_{t+1}/D_t$ 衡量的是跨越“平均滞后 $+1$”周的增长，取对数后 $\ln G\approx\sum_{g}\ln R_g$，求和跨越若干世代。各州各周的增长趋势存在异质性，若每一世代的对数增长带有一份与规模无关的离散，跨越的世代越多，累积离散越大。这部分离散不随 $D$ 衰减，于是全部进入截距。原文表 1 的阶段平均增长比随平均滞后同向变化（单周 1.4802、主权重 1.6706、等权 1.9929），与跨度越长累积增长越大的方向一致。原文第 508、665 行还指出截距包含“滞后卷积平滑失配”等代理偏差，这同样与权重有关。

因此 $a$ 应解读为“周增长比的规模无关离散”，不宜等同于 $v_R/\bar R^2$。修改时建议在原文第 446 行后写明 $G$ 与 $R_t$ 的口径差，并把三种权重下的 $a$ 作为证据报告。

\subsubsection{原文表 1 的读法}

原文表 1（第 447--464 行）对三种滞后权重分别重新筛选 $D_{j,t}\ge5$ 的样本并重做 FGLS，表~\ref{tab:em-lag} 摘录其数值并补上两列权重特征。

\begin{table}[!tbp]
\centering
\small
\caption{原文表 1 摘录：不同滞后权重下的仿射模型参数及权重特征}
\label{tab:em-lag}
\setlength{\tabcolsep}{3pt}
\begin{tabular}{L{2.6cm}L{1.0cm}L{1.2cm}L{1.1cm}L{1.3cm}L{1.2cm}L{2.6cm}L{1.1cm}L{1.1cm}}
\toprule
权重 & $\sum w^2$ & 平均滞后（周） & $N$ & $\bar R_{\mathrm{early}}$ & $b$ & 州聚类 95\% 区间 & $a$ & $m_\times$ \\
\midrule
主权重 $(0.65,0.25,0.10)$ & 0.495 & 1.45 & 1267 & 1.6706 & 2.2434 & $[1.5019, 3.0161]$ & 0.0991 & 22.65 \\
三周等权 & 1/3 & 2.0 & 1240 & 1.9929 & 1.9262 & $[1.0706, 2.8074]$ & 0.1524 & 12.64 \\
单周 $(1,0,0)$ & 1 & 1.0 & 1290 & 1.4802 & 3.1337 & $[2.2031, 4.0668]$ & 0.0682 & 45.96 \\
\bottomrule
\end{tabular}
\end{table}

$N$ 列不同，原因是每种权重重新计算 $D_{j,t}$ 并重新筛选；$\bar R_{\mathrm{early}}$ 列不同，原因是 $G=C/D$ 本身依赖权重。斜率列对应式~\ref{eq:em-delta}，随 $\sum w^2$ 增大；截距列对应本小节的口径差机制，随平均滞后增大；最后一列 $m_\times=b/a$ 同时受两者影响。平均滞后越短，$\sum w^2$ 越大，$b$ 越大而 $a$ 越小，两种机制把 $b/a$ 推向同一方向，于是 $m_\times$ 从等权的 12.64 经主权重的 22.65 升到单周的 45.96。原文把这一敏感性解读为“经验拐点应作为系统相关的量级参照”，本节的两个机制给出了它的来源。三种规格的斜率区间都远离零点，规模项为正的定性结论对权重稳健。

\subsection{FGLS 的一致性与权重}\label{sec:em-fgls}

原文第 508 行说明估计采用一阶可行广义最小二乘，权重取拟合值平方的倒数 $1/\hat\mu^2$。记回归元 $x=(1,1/D)^\top$，参数 $\beta=(a,b)^\top$，权重 $W(D)>0$ 只依赖 $D$，加权最小二乘的估计方程为
\begin{equation}\label{eq:em-ee}
\sum_i W(D_i)\,x_i\,(L_i-x_i^\top\beta)=0 .
\end{equation}
\begin{enumerate}
\item 由塔式性质，在真值 $\beta_0$ 处有 $\mathbb{E}[W(D)x(L-x^\top\beta_0)]=\mathbb{E}\big[W(D)x\,\mathbb{E}[L-x^\top\beta_0\mid D]\big]=0$，所以只要均值设定 $a+b/D$ 正确，任意 $D$ 的正函数作权重都使式~\ref{eq:em-ee} 给出一致估计 \citep{wooldridge2010econometric}。
\item 第一步 OLS 在同一条件下一致，于是拟合值 $\hat\mu$ 一致，由连续映射定理，以 $1/\hat\mu^2$ 为权重的一步 FGLS 也一致。
\item 由 Gauss--Markov 型论证，最优权重是 $1/\operatorname{Var}(L\mid D)$；把损失写成 $L=\mu(D)\varepsilon$，$\operatorname{Var}(\varepsilon\mid D)=\mathrm{CV}^2(D)$，则 $\operatorname{Var}(L\mid D)=\mu^2(D)\mathrm{CV}^2(D)$，只有 $\mathrm{CV}^2$ 为常数时 $1/\hat\mu^2$ 才是最优权重。
\item 分箱计算的 $\operatorname{Var}(L)/\mathbb{E}[L]^2$ 就是 $\mathrm{CV}^2$ 的箱内估计，它从 $[5,20)$ 箱的 5.32 降到 $[600,5000)$ 箱的 0.99，说明 $\mathrm{CV}^2$ 随规模下降，恰好违背第 3 步的条件。
\end{enumerate}
原文第 512 行称这一递减“为 FGLS 采用 $1/\hat\mu^2$ 权重提供了依据”，逻辑方向反了（I-24）。不过权重只影响效率，不影响一致性。蒙特卡洛核验在实证参数附近做了 60 次重复：$\mathrm{CV}^2$ 为常数时 $\hat a$ 的样本平均为 0.1003（SE 0.0009），$\hat b$ 为 2.1892（SE 0.0325）；$\mathrm{CV}^2$ 递减时分别为 0.1007（SE 0.0011）与 2.1895（SE 0.0451）。真值为 0.0991 与 2.2434，$\hat b$ 约低 2.4\%，在 1.2--1.7 个 SE 之内。$\hat b/\hat a$ 的样本平均为 21.93 与 22.00、中位数为 22.00 与 21.27，真值 22.64，没有出现比值正偏。修改建议有两种：改用 $1/(\hat\mu^2\widehat{\mathrm{CV}}{}^2(D))$ 权重，其中 $\widehat{\mathrm{CV}}{}^2(D)$ 对 $\log D$ 作线性拟合；或者直接写明“均值设定正确时任意 $D$ 的函数作权重均给出一致估计，权重仅影响效率”，并给出 OLS 对照。

\subsection{州聚类 Bootstrap}

原文的区间由 1,500 次按州聚类的 Bootstrap 给出。每次以州为单位有放回地重抽 51 个州，保留每个州内的全部周，按第~\ref{sec:em-fgls} 小节的估计方程重新拟合 FGLS，取百分位区间，$b/a$ 的区间取比值的百分位。聚类 Bootstrap 的有效性依赖三个条件：簇数趋于无穷，簇间独立，簇内允许任意相关 \citep{cameron2008bootstrap,cameron2015practitioner}。51 个簇处在可用范围的下限附近；同一周跨州的共同冲击（例如全国性的节假日填报积压）会破坏簇间独立。原文做了两种单向敏感性重抽：在每次重抽中重估样本平均基准，$b$ 的区间为 $[1.5170, 2.9807]$；按观测周聚类，$b$ 的区间为 $[1.4728, 3.1191]$。两者都只处理一个维度，更完整的做法是按州与周双向聚类 \citep{cameron2011robust}。

蒙特卡洛核验构造 51 个州 $\times$ 25 周的数据，每次 Bootstrap 重抽 99 次，重复 36 次。95\% 区间对 $b$ 与 $m_\times$ 的覆盖率都是 0.944（蒙特卡洛 SE 约 0.036），没有明显欠覆盖，但重复次数有限，只能排除严重欠覆盖。$m_\times$ 区间的相对宽度（上界减下界再除以点估计）在仿真中的中位数为 0.78（四分位 0.65--0.85），原文 $[13.7, 33.7]/22.6$ 为 0.88，处在同一量级。原文“1,500 次重抽中无一次出现 $b\le0$”等价于 Bootstrap 单侧 $p<1/1500$，这是对“规模项为正”最直接的推断表述。

\subsection{滞后工具变量}

原文第 515 行用滞后 3--5 周的规模 $D_t^{\mathrm{iv}}$ 作工具做两阶段最小二乘。全初期汇集样本 $N=1{,}051$，以 $\bar R_{\mathrm{early}}=1.6706$ 归一化时第二阶段 $b=2.5305$（$a=0.0812$），对应 OLS 为 1.3525；以子样本自身平均 1.6266 归一化时 $b=2.8501$（$a=0.0807$），OLS 为 1.4581。第一阶段未聚类 $F=1147.1$，按州聚类稳健 $F=312.9$；分季运行时稳健 $F$ 为 94.5、119.4、76.8，$b_{\mathrm{iv}}$ 为 3.2787、2.1867、3.6131，均高于对应的 OLS。汇集样本中差值 $b_{\mathrm{iv}}-b_{\mathrm{ols}}=1.3921$，区间 $[0.43, 2.41]$。

$b_{\mathrm{iv}}>b_{\mathrm{ols}}$ 可以用经典测量误差模型解释 \citep{fuller1987measurement,carroll2006measurement}。若观测回归元 $x=x^*+u$，测量误差 $u$ 与真值 $x^*$ 及回归误差独立，则 OLS 斜率的概率极限为
\begin{equation*}
\operatorname{plim}\hat b_{\mathrm{ols}}=b\cdot\frac{\operatorname{Var}(x^*)}{\operatorname{Var}(x^*)+\operatorname{Var}(u)},
\end{equation*}
向零衰减。滞后 3--5 周的规模与当期测量误差独立，用它作工具可以消除衰减。两种偏差的作用位置不同。第~\ref{sec:em-delta} 小节的分母噪声进入的是被解释变量 $L$ 本身，其幅度与真实的 $1/(\rho\Lambda)$ 成正比，作用方向是抬高 $b$；工具变量用滞后规模预测真实规模，这部分抬高会被保留下来。两种偏差方向相反，工具变量只纠正其中一种。

工具变量的排除限制更难成立（I-19）。滞后 3--5 周的规模与当前所处的流行阶段相关，而阶段会影响式 28 中的偏差项 $(\mu/D-\bar R)^2$，所以 $\operatorname{Cov}(Z,e)\ne0$ 的可能无法排除。第一阶段 $F$ 很大，只说明工具的相关性足够，不能检验外生性。原文自己的定位“在替代性滞后测度下，规模项为正的形态与显著性依然保持稳健”是恰当的，中期报告中“排除内生性”的说法过强，答辩时宜沿用原文的定位。

\subsection{误差预算与大规模端}\label{sec:em-budget}

原文第 554 行把初期波段的全体平均相对平方误差 0.181 分解为截距份额 54.6\%、规模项份额 45.1\%、模型外残差不足 0.3\%。分解的代数依据是对原文式 31 两端取样本平均：
\begin{equation}\label{eq:em-budget}
\overline{L}=a+b\,\overline{1/D}+\overline{e},\qquad
\frac{b\,\overline{1/D}}{a}=\frac{m_\times^{\mathrm{emp}}}{D_{\mathrm{harm}}},\quad D_{\mathrm{harm}}=\big(\overline{1/D}\big)^{-1}.
\end{equation}
截距份额为 $0.0991/0.1813=54.6\%$（分母用未取整的 0.1813）。由规模项份额 45.1\% 反推 $\overline{1/D}=0.451\times0.181/2.2434=0.0365$，样本的调和平均规模 $D_{\mathrm{harm}}\approx27.4$。式~\ref{eq:em-budget} 右式给出份额之比 $22.6/27.4=0.826$，与 $45.1/54.6$ 相等。这一关系说明误差预算由“经验交叉点相对样本调和平均规模的位置”决定：样本调和平均规模略高于 22.6，所以截距份额略过半。

FGLS 带权重，残差的样本平均不必为零，所以三项之和不是 100\%。原文第 551 行的“回顾口径 54.8\%/45.2\%”是去掉残差后重新归一的份额，$54.6/(54.6+45.1)=54.8\%$；同一规格在第 551 行与第 554 行用了两种归一化，正文没有说明（I-18），修改时宜统一口径或加注。第 551 行另报告三种事前滚动规格下的份额：全部历史扩张窗 54.9\%/45.1\%，26 周滚动池 39.1\%/60.9\%，逐州扩张窗 41.3\%/58.7\%；达峰期条件方差口径下为 46.7\% 对 55.6\%，两项合计 102.3\%，残差份额 $-2.3\%$。原文据此得出“约四至五成半（39\%--55\%）的误差份额与规模无关”。

大规模端的对照是一个机械推论。$D\ge250$ 组共 244 个观测，模型预测为 $a+b\,\mathbb{E}[1/D\mid D\ge250]\le0.0991+2.2434/250=0.108$，实测为 0.104，与截距接近是拟合模型本身的结果；该组与全体之比为 0.574，即原文所说“约六成”。原文已正确说明这一一致性“不构成对其机制来源的识别”。

\subsection{空间汇聚与覆盖率冲击}\label{sec:em-space}

\subsubsection{空间汇聚（I-06、I-20）}

原文第 617 行把 51 个辖区按周加总，全美序列初期波段的平均周规模约 6662 例，相对平方误差降到 0.0348（中位数 0.0115，$N=28$），约为单州平台 0.10 的三分之一。原文用空间相关解释这一降幅，给出的汇总环境方差如下式所示：
\begin{equation*}
v_R^{\mathrm{agg}} = v_R \left[\bar{c} + \frac{1-\bar{c}}{J}\right]
\tag*{原文第 617 行}
\end{equation*}
该式的推导如下。记汇总增长 $R_{\mathrm{agg}}=\sum_j\pi_jR_j$，各辖区环境方差同为 $v_R$、两两相关系数为 $c_{jl}$，由方差的双线性展开得 $\operatorname{Var}(R_{\mathrm{agg}})=v_R\big[\sum_j\pi_j^2+\sum_{j\ne l}\pi_j\pi_lc_{jl}\big]$；取等权 $\pi_j=1/J$ 与共同相关 $\bar c$，括号内为 $1/J+(1-1/J)\bar c=\bar c+(1-\bar c)/J$。权重不等时，$J$ 换成有效个数 $1/\sum_j\pi_j^2$，全国按人口加总时有效个数小于 51。

原文没有报告隐含的 $\bar c$，可以反推。全国截距约为 $0.0348-2.2434/6662=0.0345$；若州级截距全部是环境项，则 $\bar c+(1-\bar c)/51=0.0345/0.0991=0.348$，解得 $\bar c\approx0.34$。这个数只是粗略指示：$N=28$，中位数 0.0115 远低于样本平均 0.0348，分布偏斜；州级截距含有第~\ref{sec:em-intercept} 小节的非环境成分。原文第 168 行称环境方差“无法随横截面聚合消除”，与第 617 行的汇总结论矛盾（I-06）；准确的表述是 $v_R$ 无法通过同一环境内扩大规模消除，跨相关性不完全的环境聚合可以削减其 $(1-\bar c)$ 部分。

\subsubsection{2024 年覆盖率骤降（I-21）}

原文第 613 行报告：第一阶段观测增长率标准差 0.624、偏度 3.29（法定强制子段为 0.620 与 3.45）；第二阶段标准差升到 1.823（扩张 192.1\%）、偏度 14.98；第三阶段回落到 0.649 与 2.68。四分位距（0.638 到 0.594）与中位绝对离差（0.298 到 0.299）基本不变，扩张集中在右尾。

均匀漏报能解释多少扩张，可以给出上界。设覆盖率下降等价于有效报告概率从 $\rho_1$ 均匀二项抽稀到 $\rho_2$。在命题 3 的 $V(n,\rho)=v_R+\bar R/(\rho n)+(\bar R^2+v_R)/(kn)$ 中，只有 $\bar R/(\rho n)$ 一项随 $\rho$ 变化，恰好放大 $\rho_1/\rho_2$ 倍，其余两项不变；分母噪声项 $(1-\rho)/\rho$ 在有效报告概率只有百分之一量级时 $(1-\rho)\approx1$，也近似放大 $\rho_1/\rho_2$ 倍。于是方差至多放大 $\rho_1/\rho_2$ 倍：
\begin{equation*}
\frac{\mathrm{SD}_2}{\mathrm{SD}_1}\le\sqrt{\rho_1/\rho_2}=\sqrt{0.911/0.546}=1.29 .
\end{equation*}
实测 $1.823/0.624=2.92$，远超上界，均匀漏报无法解释。这一定量结果支持原文“结构性非随机扰动”的判断（医院进出填报体系），与命题 3 的联系只能是定性的。另外，命题 3 的非对称性在住院情景下本来就很弱：同比例增加 $m$ 时沿 $n$ 与沿 $\rho$ 的边际降幅之比为 $A/\bar R=1+\rho(\bar R^2+v_R)/(k\bar R)$，在 $\rho\le0.021$ 时只有 1.019--1.180（I-05）。

\subsection{达峰期与稳健性表}

\subsubsection{达峰期偏差分解与尺度等变}

原文把各季全国达峰周（2022-12-03、2023-12-30、2025-02-08）及其后 7 周定义为达峰与回落期，共 1220 个州–周，$\bar R_{\mathrm{peak}}=0.8075$（标准差 0.3048）。以初期基准预测时，总误差含一个与规模无关的基准偏移项，原文的写法如下：
\begin{equation*}
\left(\bar{R}_{\mathrm{early}}-\bar{R}_{\mathrm{phase}}\right)^2/\bar{R}_{\mathrm{early}}^2
\tag*{原文第 558 行}
\end{equation*}
它来自平方展开 $(G-\bar R_{\mathrm{early}})^2=(G-\bar R_{\mathrm{peak}})^2+2(G-\bar R_{\mathrm{peak}})(\bar R_{\mathrm{peak}}-\bar R_{\mathrm{early}})+(\bar R_{\mathrm{peak}}-\bar R_{\mathrm{early}})^2$；$\bar R_{\mathrm{peak}}$ 是阶段样本平均，交叉项在全阶段样本上取平均后为零。代入得 $(0.8075-1.6706)^2/1.6706^2=0.74494/2.79090=0.2669$，原文称其占 88.9\%，反推总损失约 0.300，落在原文表 4 的 0.292--0.312 之内。按各箱自身平均计算的偏差占比从 59.5\% 升到 94.8\%，总体份额不代表每个箱。

阶段间比较斜率时，需要用到 FGLS 的尺度等变性：若 $L'=cL$，OLS 系数乘以 $c$，拟合值 $\hat\mu'=c\hat\mu$，权重整体乘以 $c^{-2}$，而加权最小二乘的解对权重的公共倍数不变，所以 FGLS 系数也乘以 $c$。据此，达峰期以自身平均归一化的 $b=7.9272$ 换到共同分母 $\bar R_{\mathrm{early}}^2$ 下为 $7.9272\times(0.8075/1.6706)^2=7.9272\times0.23364=1.8521$；未归一化口径下初期为 $2.2434\times2.7909=6.2611$，达峰期为 $7.9272\times0.65206=5.1690$。

\subsubsection{原文表 3（竞争基准）的读法}

原文表 3（第 525--548 行）按规模箱比较四种基准的事前单步相对平方误差，表~\ref{tab:em-baselines} 摘录了其中五行。

\begin{table}[!tbp]
\centering
\small
\caption{原文表 3 摘录：初期波段四种预测基准的相对平方误差}
\label{tab:em-baselines}
\begin{tabular}{L{2.6cm}L{1.4cm}L{2.2cm}L{2.2cm}L{2.2cm}L{2.2cm}}
\toprule
规模箱 $D_t$ & 样本数 $N$ & 阶段平均基准 & Cori 2 周窗 & Cori 3 周窗 & 朴素平移 \\
\midrule
$[5, 20)$ & 319 & 0.3292 & 0.6576 & 0.5606 & 1.2563 \\
$[100, 250)$ & 220 & 0.1118 & 0.0903 & 0.0870 & 0.3999 \\
$[600, 5000)$ & 92 & 0.1054 & 0.0983 & 0.0982 & 0.1360 \\
全体样本平均 & 1267 & 0.1813 & 0.2876 & 0.2490 & 0.6367 \\
全体样本中位数 & 1267 & 0.0653 & 0.0658 & 0.0626 & 0.1732 \\
\bottomrule
\end{tabular}
\par\smallskip
\begin{minipage}{0.92\linewidth}\footnotesize
注：Cori 两列因局部比值非正分别剔除 23 与 16 个点，有效样本为 1,244 与 1,251。各列以各自预测值归一化。
\end{minipage}
\end{table}

第三列与原文表 2 的第三列相同。Cori 基线 \citep{cori2013new} 用短窗口估计的 $\widehat R_t$ 替代 $\bar R$，相当于把理论中的已知基准换成一个自身带有抽样方差的估计量。由命题 1，这个估计量的方差随规模按 $1/m$ 衰减，于是它叠加在规模项上。两条 Cori 基线的仿射斜率达到 6.0637 与 4.9836，交叉点为 74.7 与 65.3，小规模箱误差高于阶段平均基准。大规模箱中 Cori 基线（约 0.087--0.098）略低于阶段平均基准（0.103--0.112），对应原文式 28 的偏差项被时变参数部分吸收。四列各自以预测值归一化，不能直接横向比较。原文第 521 行在 1,244 个共同样本上改用回顾性阶段平均作共同分母，Cori 2 周窗与朴素平移的整体误差变为 0.2230 与 0.2738，小规模箱中阶段平均基准（0.3275）优于时变基准（0.4149）与朴素预测（0.4951）。原文第 551 行的事前滚动核验同属这一类稳健性检验：全部历史扩张窗、26 周滚动池、逐州扩张窗的 $b$ 为 9.1579、9.4562、13.2666，$b/a$ 为 22.5、42.7、38.9，截距升到 0.4063、0.2217、0.3409。基准越不准，截距越高，与式 28 的偏差项方向一致。

\subsubsection{原文表 4（达峰期分箱）的读法}

原文表 4（第 562--580 行）用初期基准 $\bar R_{\mathrm{early}}$ 计算达峰期各箱的相对平方误差，表~\ref{tab:em-peak} 照录全部六行。

\begin{table}[!tbp]
\centering
\small
\caption{原文表 4 照录：以初期基准预测达峰与回落期时的分箱误差}
\label{tab:em-peak}
\begin{tabular}{L{2.8cm}L{2.0cm}L{2.8cm}L{2.2cm}}
\toprule
规模箱 $D_t$ & 样本数 $N$ & 相对均方误差 & 中位数 \\
\midrule
$[5, 20)$ & 33 & 0.2998 & 0.2361 \\
$[20, 50)$ & 140 & 0.2927 & 0.2591 \\
$[50, 100)$ & 162 & 0.3120 & 0.3085 \\
$[100, 250)$ & 312 & 0.3049 & 0.3005 \\
$[250, 600)$ & 311 & 0.2998 & 0.2976 \\
$[600, 5000)$ & 262 & 0.2916 & 0.2987 \\
\bottomrule
\end{tabular}
\end{table}

这张表对应上文的偏差分解。各箱误差在 0.292--0.312 之间几乎持平，因为 0.2669 的常数偏移项与规模无关且占据主导。中位数与箱内样本平均接近（例如 $[250,600)$ 箱 0.2976 对 0.2998），与初期波段中位数远低于样本平均的重尾形态（表~\ref{tab:em-scale}）不同，这也符合“一个共同的常数偏移支配损失”的解读。改用 $\bar R_{\mathrm{peak}}$ 作基准后，仿射拟合给出 $a=0.0665$ $[0.0467, 0.0897]$、$b=7.9272$ $[5.4239, 11.2248]$，$b/a=119.2$，规模项依然为正。

\subsubsection{原文表 5（阶段规模）的读法}

原文表 5（第 588--603 行）把两阶段的 $b$ 放在三种损失口径下比较，表~\ref{tab:em-stage} 照录并补上换算因子。

\begin{table}[!tbp]
\centering
\small
\caption{原文表 5 照录：两阶段规模项 $b$ 在三种损失口径下的比较及换算关系}
\label{tab:em-stage}
\begin{tabular}{L{4.2cm}L{1.6cm}L{1.8cm}L{5.4cm}}
\toprule
损失口径 & 初期 $b$ & 达峰期 $b$ & 换算关系 \\
\midrule
未归一化 $(G-\bar R_{\mathrm{phase}})^2$ & 6.2611 & 5.1690 & $2.2434\times2.7909$；$7.9272\times0.65206$ \\
共同分母 $\bar R_{\mathrm{early}}^2$ 归一化 & 2.2434 & 1.8521 & 达峰期 $5.1690/2.7909$ \\
各阶段自身平均归一化 & 2.2434 & 7.9272 & 直接拟合值 \\
\bottomrule
\end{tabular}
\end{table}

三行之间只差一个常数因子，依据是上文的 FGLS 尺度等变性。第三行显示达峰期 $b$ 更大，那只是分母 $0.8075^2$ 远小于 $1.6706^2$ 的结果；在共同分母与未归一化口径下，达峰期的 $b$ 都低于初期。原文在同一样本、共同分母下估计了四种规格的斜率差：分别 OLS 为 2.3049 对 1.3866，差 $-0.9183$；联合 OLS 交互项同为 $-0.9183$，州聚类区间 $[-2.33, 0.59]$；分别 FGLS 差 $-0.3913$，联合 FGLS 交互项相同。带完全交互的联合模型中，交互项系数等于分别拟合的斜率差，这是代数恒等。共同支撑集 $[50,600)$ 上原文报告共同分母口径为 1.6326 对 2.9264，结论相同：点估计为负，区间包含零，数据不支持达峰期规模项更强，也不足以断言其更弱。

\subsubsection{原文表 6（扩展验证）的读法}

原文表 6（第 629--650 行）把同一个仿射模型用到两组外部数据上，表~\ref{tab:em-extended} 摘录每组的主规格。

\begin{table}[!tbp]
\centering
\small
\caption{原文表 6 摘录：外部数据上的仿射模型主规格}
\label{tab:em-extended}
\setlength{\tabcolsep}{4pt}
\begin{tabular}{L{3.4cm}L{3.0cm}L{0.9cm}L{1.3cm}L{2.7cm}L{1.1cm}L{1.2cm}}
\toprule
数据 & 筛选口径 & $N$ & $b$ & 聚类 95\% 区间 & $a$ & $m_\times$ \\
\midrule
FluSight 集成预测 2022--23 & 州级事前筛选 $\widehat Y\ge5$ & 1177 & 7.8045 & $[5.3582, 10.5260]$ & 0.1295 & 60.28 \\
欧洲 11 国 COVID-19 & 事前规模 $D_t\ge5$、$g>0$ & 404 & 33.9223 & $[14.6021, 74.9873]$ & 0.2721 & 124.7 \\
\bottomrule
\end{tabular}
\end{table}

FluSight 一组 \citep{mathis2024evaluation} 的自变量是集成模型发布的预期住院规模 $\widehat Y_{j,t}$，被解释变量是集成预测的相对平方误差，所以它检验的是“误差随规模按 $a+b/D$ 衰减”的形态在外部模型上是否出现，不对应理论中的 $v_R$ 与 $A$。原文报告的另一个关键数字是达峰期偏差占比：集成模型下只有 15.5\%（上升期 32.3\%，回落期 1.5\%），对比单代常数基准的 88.9\%，说明自适应模型能大幅压缩式 28 的偏差项。欧洲 COVID-19 一组 \citep{flaxman2020estimating} 使用 7 天平滑的日度确诊，规模单位是“例/天”，增长期样本与截尾子样本的 $b$ 为 18.5337 与 12.7632，交叉点为 88.1 与 67.7。两组数据的指标定义、时间分辨率与筛选规则都与 NHSN 不同，所以各行的 $m_\times$ 不能互相比较，原文第 627 行也写明“不能把这些拐点视为跨病种通用阈值”。这张表能支持的结论只有一条：规模项为正的形态在不同模型与病种中重复出现。
